\chapter{Diagonalization and the First Incompleteness Theorem}
\label{ch:first-incompleteness}

\section{The diagonal lemma}
\label{sec:first-incompleteness-01}

Let \(\mathsf{diag}(n)\) be the code of the result of substituting the numeral
\(\bar n\) for the distinguished free variable in the one-variable formula
with code \(n\).  On malformed codes or codes with other free variables,
use the fixed sentence code \(\ulcorner0=0\urcorner\), as specified in
\cref{def:p05-rep-diagonal-default}.  The
function \(\mathsf{diag}\) is primitive recursive.
In the construction below, the displayed free variable \(x\) denotes
the distinguished variable \(x_0\).  A formula presented with another
free variable is first capture-avoidingly renamed to use \(x_0\), and
all auxiliary bound variables are chosen fresh before coding.

\begin{theorem}[Diagonal lemma]
\label{thm:diagonal-lemma}
Let \(T\) be any theory extending \(\mathsf Q\).  For every
\(\Lang_A\)-formula \(\theta(x)\) with at most the displayed free variable,
there is a sentence \(G\) such that
\[
  T\derives G\leftrightarrow
  \theta(\overline{\ulcorner G\urcorner}).
\]
\textit{Status: complete proof.}  The full coding and graph-based
derivations are in \cref{annex:p05-rep-diagonal-graph,
annex:p5.inc.diagonal-coding,annex:p5.inc.diagonal}.
\end{theorem}

See also \cite{BoolosBurgessJeffrey2007}.  The following
construction explains the fixed point used by subsequent complete proofs.
Let \(D(x,y)\) represent the function \(\mathsf{diag}\) in \(\mathsf Q\), so
for each standard \(n\),
\[
  \mathsf Q\derives\forall y\,
  \bigl(D(\bar n,y)\leftrightarrow
  y=\overline{\mathsf{diag}(n)}\bigr).
\]
Define
\[
  \beta(x)\defeq\exists y\,(D(x,y)\land\theta(y)).
\]
Let \(b=\ulcorner\beta\urcorner\), and let
\(G\) be the sentence \(\beta(\bar b)\).  By construction,
\(\mathsf{diag}(b)=\ulcorner G\urcorner\).  Representability of
\(\mathsf{diag}\) gives, already in \(\mathsf Q\),
\[
\begin{aligned}
  G
  &\leftrightarrow \exists y\,(D(\bar b,y)\land\theta(y))\\
  &\leftrightarrow
  \theta(\overline{\mathsf{diag}(b)})\\
  &\leftrightarrow\theta(\overline{\ulcorner G\urcorner}).
\end{aligned}
\]
Since \(T\supseteq\mathsf Q\), \(T\) proves the same biconditional.
The explicit graph in \cref{thm:p05-rep-diagonal-graph} supplies both
existence and unrestricted-output uniqueness at the standard code used
here; \cref{thm:p05-inc-diagonal-coding,thm:p05-inc-diagonal} give the
complete coding and inference derivations.  Numeralwise output instances
alone would not suffice for the displayed universal-output graph clause.

\begin{remark}
The lemma does not create a sentence that literally contains its own
inscription.  It creates a sentence provably equivalent to the result of
applying \(\theta\) to the numeral for that sentence's code.
\end{remark}
